% ================================================================
% chapters/ch2_literature.tex
% Chapter 2 — Literature Review
% ================================================================

\chapter{Literature Review}
\label{ch:literature}

This chapter reviews the body of knowledge relevant to this
thesis across six areas: the global CBDC landscape, privacy
in digital payment systems, cryptographic schemes for CBDC
privacy, known de-anonymization attacks, existing privacy
measurement frameworks, and Hyperledger Fabric as the
reference deployment platform. The chapter concludes by
identifying the specific research gap that this thesis fills.

% ================================================================
\section{Central Bank Digital Currency --- Global Overview}
\label{sec:cbdc_global}
% ================================================================

A Central Bank Digital Currency is a digital form of a
country's sovereign currency, issued and backed directly by
the central bank, representing a liability of the central bank
rather than of a commercial intermediary
\cite{mancinigrifoli2018}. This distinguishes a CBDC from
commercial bank deposits, from privately issued stablecoins,
and from decentralised cryptocurrencies such as Bitcoin, none
of which carry the direct backing of a monetary authority.

The taxonomy of CBDCs is most usefully divided along two
dimensions. The first is the target user: a \textit{retail}
CBDC is designed for use by households and businesses in
everyday transactions, while a \textit{wholesale} CBDC is
restricted to financial institutions for interbank settlement.
This thesis is concerned exclusively with retail CBDC, which
is the design most relevant to financial inclusion objectives
and to the privacy risks analysed here. The second dimension
is the structural model: a \textit{token-based} CBDC assigns
value to a cryptographic token that the holder controls
through a private key, while an \textit{account-based} CBDC
records balances against verified identities in a centralised
ledger \cite{biscgide2024}. This thesis models token-based
CBDC, in which the cryptographic design of the token directly
determines the privacy properties of each transaction.

Globally, the pace of CBDC exploration has accelerated
significantly. The Bank for International Settlements reported
that over 130 countries were actively researching or piloting
CBDCs by 2024, representing more than 98\% of global GDP
\cite{bis2020cbdc}. Among operational deployments, China's
digital yuan (e-CNY) is the most advanced, having processed
billions of yuan across multiple major cities. The Bahamas
Sand Dollar became the world's first fully deployed retail
CBDC, followed by Jamaica's JAM-DEX and Nigeria's eNaira.
In Europe, the digital euro project is advancing through
its preparation phase, with privacy cited as a core design
requirement by the European Central Bank.

Morales-Resendiz et al. \cite{moralesresendiz2021}
provide a detailed peer review of three retail CBDC pilot
programs --- the Bahamas Sand Dollar, Uruguay's e-Peso, and
Sweden's e-Krona --- drawing common lessons on design,
operation, and implementation. Their analysis highlights that
despite differences in motivation and technical architecture,
all three pilots converge on the importance of a hybrid
operational model in which the central bank retains core
infrastructure while the private sector handles user-facing
services. The privacy implications of this hybrid model are
directly relevant to this thesis: under hybrid architecture,
transaction metadata is visible to multiple participants
beyond the central bank alone.

For developing and emerging economies specifically,
Ozili \cite{ozili2022} presents a balanced analysis of
the financial inclusion arguments for and against CBDC
adoption. The arguments in favour centre on CBDC's ability
to extend digital financial services to unbanked populations,
reduce transaction costs, and function in offline
environments with limited connectivity. The arguments against
emphasise the cost of digital devices for low-income users,
the risk that identification and KYC requirements create
new barriers, and the possibility that CBDC design priorities
may not align with financial inclusion objectives. Nepal's
context sits squarely within this debate: the financial
inclusion opportunity is real, but so is the risk that a
poorly designed CBDC could expose the privacy of the very
populations it is intended to serve.

The BIS CGIDE \cite{biscgide2024} proposes a reference
architecture for retail CBDC that separates transaction data
from identity information, with identity retained by private
intermediaries. This separation is presented as a privacy
safeguard. However, the proposal also acknowledges that this
separation does not fully resolve the metadata privacy
problem: transaction data that does not contain identity
fields may still be sufficient to re-identify users through
pattern analysis, which is the precise problem this thesis
quantifies.

% ================================================================
\section{Privacy in Digital Payment Systems}
\label{sec:privacy_payments}
% ================================================================

Privacy in the context of digital payments is not a single
concept but a spectrum of distinct properties that must be
carefully distinguished.

\textit{Anonymity} means that a transaction cannot be linked
to any real-world identity, even by the payment system
operator. Physical cash is the canonical example: the central
bank has no record of who holds which banknote or what it
was used to purchase.

\textit{Pseudonymity} means that transactions are linked to
a persistent identifier --- a pseudonym --- rather than a
real name. The system operator may or may not know the
mapping between pseudonym and real identity. Hyperledger
Fabric operates on this model: each user has an X.509
certificate-based identity, and the central bank knows the
mapping, but transactions on the public channel appear under
wallet addresses rather than names.

\textit{Confidentiality} means that transaction contents are
hidden from third parties, even if the parties to the
transaction are identifiable. This is the property provided
by encrypted messaging and by ZKP range proofs that hide
transaction amounts.

The critical insight motivating this thesis is that
pseudonymity without metadata privacy is insufficient. Garratt
and van Oordt \cite{garratt2020} argue that privacy in
digital payments should be treated as a public good --- one
that market mechanisms alone will not provide at the socially
optimal level --- and that central bank design choices have
a direct role in determining how much metadata privacy
citizens receive. They distinguish between operational
privacy, which concerns who can see transaction data, and
informational privacy, which concerns what can be inferred
from transaction data even by parties with legitimate access.
This thesis addresses informational privacy: the risk of
inference from metadata by parties who have legitimate access
to the ledger but should not be able to re-identify users
from it.

The metadata problem arises because even when names are
removed from transaction records, the combination of
observable fields --- amount, time, merchant, wallet
address, frequency --- creates a unique behavioural
fingerprint. A user who pays NPR 4,738 at a specific
merchant at 2:47am in a hill rural zone has produced a
transaction record that may be unique to exactly one person
in the population. No cryptographic attack is required. The
information is present in the metadata structure itself.
This phenomenon is well documented in the digital payments
literature and has been demonstrated empirically in the
blockchain context by Koshy et al. \cite{koshy2014},
who showed that Bitcoin users could be identified from
transaction patterns observable in peer-to-peer network
traffic, even without any cryptographic attack on Bitcoin's
elliptic curve signature scheme.

% ================================================================
\section{Cryptographic Schemes for CBDC Privacy}
\label{sec:crypto_schemes}
% ================================================================

The choice of cryptographic scheme in a token-based CBDC
determines what metadata is observable to an adversary with
ledger access. Six schemes are relevant to this thesis,
spanning from the weakest to the strongest metadata privacy
guarantees.

\textbf{Plain Hash Token.} The simplest token design uses a
cryptographic hash of the token serial number as the token
identifier. SHA-256 is the standard hash function used in
practice. This design provides no metadata hiding: all
transaction fields --- amount, timestamp, merchant, wallet
address --- are recorded exactly on the ledger. It is
included in this thesis as the baseline worst-case scheme.

\textbf{Blind Signature.} The blind signature scheme was
introduced by Chaum \cite{chaum1983} as a mechanism
for issuing unforgeable electronic cash tokens without the
bank being able to link the issued token to the spending
event. Asghar \cite{asghar2011} provides a comprehensive
survey of blind signature constructions from RSA-based to
discrete logarithm problem-based schemes. The core privacy
property is \textit{unlinkability}: the bank cannot connect
the token issuance to the token spending. In the context of
CBDC metadata, blind signatures coarsen the wallet address
and timestamp fields while leaving amount and merchant
largely unchanged.

\textbf{Ring Signature.} Ring signatures, introduced in the
context of privacy-preserving cryptocurrency by van Saberhagen
\cite{vansaberhagen2013} in the CryptoNote protocol,
allow a transaction to be signed by any member of a set of
possible signers without revealing which specific member
signed. This provides an anonymity set for the wallet field.
Noether \cite{noether2015} extended the ring signature
construction with confidential transactions to hide amounts
as well, but this amount-hiding property requires explicit
implementation. In the CBDC context, ring signatures
primarily address wallet address privacy, while amounts may
be bucketed rather than fully hidden.

\textbf{Hybrid Scheme.} The digital euro project and similar
tiered privacy proposals combine multiple cryptographic
primitives. The hybrid scheme modelled in this thesis combines
the wallet randomisation of blind signatures with the amount
bucketing of ring signature approaches, representing a
realistic intermediate option for resource-constrained
deployments such as Nepal's.

\textbf{Stealth Address.} Stealth addresses generate a
cryptographically fresh, unlinkable address for every
transaction. Unlike wallet randomisation in blind signatures,
which produces a fixed address drawn from a larger pool,
stealth addresses produce a genuinely unique address per
transaction with no mathematical link between addresses of
the same user. This eliminates wallet-based linking attacks
entirely but does not address amount or timestamp privacy.

\textbf{ZKP Token.} Zero-knowledge proof systems allow a
prover to demonstrate knowledge of a value without revealing
the value itself. Groth and Kohlweiss \cite{groth2015}
introduced one-out-of-many proofs that enable efficient
membership proofs used in CBDC privacy schemes. In a ZKP
token, all sensitive attributes are reduced to coarse
disclosure bands --- amount to a set of denominations,
timestamp to day of week, merchant to broad sector category
--- and a nullifier-based wallet provides transaction
unlinkability. ZKP tokens provide the strongest metadata
privacy of all six schemes but carry the highest
computational cost.

% ================================================================
\section{De-Anonymization Attacks}
\label{sec:deanon_attacks}
% ================================================================

The threat of metadata-based re-identification is not
theoretical. Several published attacks have demonstrated
that transaction metadata alone is sufficient to identify
individuals in supposedly anonymous datasets.

The most widely cited demonstration is the Netflix
de-anonymization study by Narayanan and Shmatikoff
\cite{narayanan2008}. Netflix released an anonymised
dataset of movie ratings for a machine learning competition,
removing all user names and replacing them with random
identifiers. Narayanan and Shmatikoff demonstrated that
knowing just eight movie ratings for a target user was
sufficient to identify that user in the dataset with high
confidence, by correlating the anonymous Netflix records with
public ratings on the Internet Movie Database. The key
finding was that the \textit{combination} of attributes ---
even individually innocuous ones --- created a unique
behavioural fingerprint. This result directly motivates the
pairwise interaction coefficient $\lambda_{ij}$ in this
thesis, which models exactly this amplification effect.

In the blockchain context, Koshy et al.
\cite{koshy2014} demonstrated that Bitcoin users could
be linked to their public addresses by analysing the timing
and structure of peer-to-peer network messages, without any
cryptographic attack. This attack is structurally analogous
to what a passive observer on a Hyperledger Fabric channel
could achieve: legitimate read access to the ledger, combined
with metadata analysis, yields re-identification without
breaking any cryptographic primitive.

For CBDC specifically, the re-identification risk is
compounded by the geographic and demographic structure of
the deployment population. In a small hill rural VDC in
Nepal with limited merchants and low transaction frequency,
the anonymity set for any individual transaction is
dramatically smaller than in a dense urban environment. The
micro-scale analysis in Chapter~\ref{ch:results} quantifies
this risk precisely.

% ================================================================
\section{Existing Privacy Measurement Frameworks}
\label{sec:privacy_frameworks}
% ================================================================

Three formal privacy measurement frameworks are directly
relevant to this thesis: k-anonymity, l-diversity, and
differential privacy. Understanding both what each framework
provides and where it falls short motivates the need for
the new metric developed here.

\textbf{k-Anonymity.} Sweeney \cite{sweeney2002}
introduced k-anonymity as a formal privacy requirement: a
dataset satisfies k-anonymity with respect to a set of
quasi-identifier attributes if every record is
indistinguishable from at least $k-1$ other records on
those attributes. A dataset with $k=1$ for some record
means that record is unique and fully identifiable. The
connection to this thesis is direct: the single-attribute
uniqueness score $s(a_i)$ in the model of this thesis
measures the fraction of users for whom $k=1$ on attribute
$a_i$. A user who is unique on an attribute --- $k=1$ ---
has zero protection from k-anonymity on that attribute.
$s(a_i)$ is therefore the k-anonymity failure rate for
attribute $a_i$.

However, k-anonymity has significant limitations in the
CBDC context. First, it considers attributes independently:
a record may satisfy k-anonymity on each attribute
individually while being uniquely identifiable from their
combination. The pairwise interaction term
$\lambda_{ij} \times s(a_i) \times s(a_j)$ in this thesis
addresses this limitation directly. Second, k-anonymity
does not provide a normalised, scheme-comparable metric.
It produces a binary assessment for each attribute rather
than a continuous risk score that can be compared across
cryptographic schemes.

\textbf{l-Diversity.} Machanavajjhala et al.
\cite{machanavajjhala2007} extended k-anonymity with
l-diversity, requiring that each equivalence class contains
at least $l$ well-represented values for sensitive
attributes. This addresses the homogeneity attack on
k-anonymous data. While l-diversity is relevant to
attribute value distribution, it does not address the
cross-attribute combination problem central to this thesis
and does not produce a normalised risk metric suitable for
scheme comparison.

\textbf{Differential Privacy.} Dwork \cite{dwork2006}
introduced differential privacy as a mathematical guarantee
that the output of a computation changes only negligibly
when any single individual is added to or removed from the
input dataset. Formally, a mechanism $M$ satisfies
$\varepsilon$-differential privacy if for any two datasets
differing in one record and any output $S$:

\begin{equation}
\Pr[M(D_1) \in S] \leq e^\varepsilon \cdot \Pr[M(D_2) \in S]
\end{equation}

The privacy budget $\varepsilon$ controls the strength of
the guarantee. Differential privacy is enforced by adding
calibrated noise to query outputs. In the CBDC context,
differential privacy on transaction amounts --- implemented
by adding noise to the recorded amount --- is one of the
mitigation strategies evaluated in Chapter~\ref{ch:results}.
The model in this thesis measures the effect of such
mitigation by applying a multiplier to $s(a_i)$, which
simulates the reduction in attribute uniqueness that noise
injection achieves.

\textbf{Limitations of Existing Frameworks.} None of the
three frameworks above directly addresses the CBDC metadata
de-anonymization problem as formulated in this thesis.
k-Anonymity and l-diversity require access to the full
dataset and do not generalise easily to the dynamic,
growing transaction ledger of a live CBDC. Differential
privacy requires implementation at the point of data
collection and does not produce a risk score comparable
across different cryptographic scheme choices. Most
importantly, none of them provides a normalised metric
bounded in $[0,1]$ with a proven theoretical maximum,
making cross-scheme comparison formally impossible under
any of these frameworks. This is the specific gap that
$R_{\mathrm{norm}}$ fills.

% ================================================================
\section{Hyperledger Fabric as CBDC Infrastructure}
\label{sec:fabric}
% ================================================================

Hyperledger Fabric is a permissioned enterprise blockchain
framework developed under the Linux Foundation's Hyperledger
project \cite{hyperledgerfabric2023}. Androulaki et al.
\cite{androulaki2018} describe Fabric's architecture in
detail: it is a modular, pluggable system supporting smart
contracts (chaincode), pluggable consensus mechanisms, and
a channel-based data partitioning model. Unlike public
permissionless blockchains such as Ethereum, Fabric requires
every participant to be identified through a Membership
Service Provider (MSP) before they can join the network.
This permissioned model makes Fabric well-suited for CBDC
deployment, where KYC and AML compliance require verified
participant identities.

The channel architecture is particularly important for
privacy: each channel is a private ledger shared only among
its members, and data on one channel is not visible to
participants on other channels. In a CBDC deployment, this
could allow different transaction tiers or different
merchant categories to be separated. However, all
participants on a channel see all transactions on that
channel, which creates the metadata exposure problem
analysed in this thesis.

Fabric transaction records contain several observable
metadata fields. Androulaki et al. \cite{androulaki2018}
identify the transaction envelope, which includes the
transaction proposal (containing chaincode ID, function
name, and arguments), the endorsement set (listing which
peers endorsed the transaction), and the block header
(containing timestamp and sequence number). In a CBDC
chaincode, the transaction arguments include the transfer
amount and recipient address. Together these fields correspond
directly to the seven attributes modelled in this thesis:
amount, timestamp, merchant (recipient organisation),
wallet (sender address), frequency (derivable from ledger
history), payment channel (derivable from chaincode ID),
and transaction zone (derivable from peer organisation
geographic configuration).

Bhattacharya et al. \cite{bhattacharya2020} study
privacy in a related Hyperledger platform (Hyperledger Indy)
and propose an attribute sensitivity scoring model for
self-sovereign identity agents. Their finding that combining
individually low-sensitivity attributes such as name and
date of birth produces highly sensitive composite
information directly parallels the $\lambda_{ij}$
interaction effect modelled in this thesis. The conceptual
connection between their attribute sensitivity framework
and the uniqueness-based approach in this thesis represents
an important cross-reference between identity privacy and
transaction metadata privacy on permissioned blockchain
platforms.

It is important to note that Fabric's permissioned identity
model does not eliminate the privacy problem this thesis
addresses. The Nepal Rastra Bank, as the network operator,
knows the real-world identity behind every wallet address
for compliance purposes. However, the metadata privacy
problem exists at a different level: it concerns what other
channel participants, insiders with partial access, and
adversaries who obtain the transaction log without the
identity mapping can infer from the observable metadata
alone. As argued in Chapter~\ref{ch:introduction}, this
is a real and consequential threat that the permissioned
identity model does not resolve.

% ================================================================
\section{Research Gap}
\label{sec:research_gap}
% ================================================================

The literature reviewed in this chapter reveals a consistent
and specific gap. Table~\ref{tab:gap} summarises the coverage
of existing work across the dimensions most relevant to this
thesis.

\begin{table}[H]
\centering
\caption{Summary of existing work and coverage gaps}
\label{tab:gap}
\begin{tabular}{lcccc}
\toprule
\textbf{Work} &
\textbf{Formal metric} &
\textbf{Multi-scheme} &
\textbf{Normalised} &
\textbf{CBDC-specific} \\
\midrule
Sweeney (2002)            & Yes & No  & No  & No  \\
Dwork (2006)              & Yes & No  & No  & No  \\
Machanavajjhala (2007)    & Yes & No  & No  & No  \\
Narayanan \& Shmatikoff (2008) & No & No & No & No \\
Koshy et al. (2014)       & No  & No  & No  & No  \\
Bhattacharya et al. (2020)& Partial & No & No & No \\
Morales-Resendiz (2021)   & No  & No  & No  & Yes \\
BIS CGIDE (2024)          & No  & No  & No  & Yes \\
\textbf{This thesis}      & \textbf{Yes} & \textbf{Yes} &
\textbf{Yes} & \textbf{Yes} \\
\bottomrule
\end{tabular}
\end{table}

The gap is clear. No existing work provides a formal,
normalised privacy risk metric that enables quantitative
comparison across multiple cryptographic scheme choices
in the specific context of token-based CBDC transaction
metadata. k-Anonymity, l-diversity, and differential
privacy provide rigorous frameworks but were not designed
for scheme comparison and do not produce normalised scores.
CBDC-specific literature provides architectural guidance
and operational lessons but no formal privacy risk
measurement tool. De-anonymization attack literature
demonstrates the threat empirically but does not generalise
to a measurement framework.

This thesis fills that gap by deriving $R_{\mathrm{norm}}$
--- a normalised metadata de-anonymization risk score with
a proven theoretical upper bound, applicable to any
token-based CBDC attribute set and cryptographic scheme
configuration --- and applying it to six schemes in Nepal's
domestic CBDC deployment context.